\ifdefined\gkver\else\def\gkver{student}\fi
\documentclass[\gkver]{gaokaozhenti}
\begin{document}

\chapter{2010年天津理综}
%% paperType: 理综物理部分

\begin{enumerate}
\item
%% number: 1
%% paperName: 2010年天津理综第 1 题 6 分
%% typeId: 1
%% type: 单选题
%% chapter: 光学
%% point: 光的电磁说
%% method: 无
%% score: 6
%% degree: 650
%% duplicateId: 0
%% body:
下列关于电磁波的说法正确的是\xzanswer{A}
\fourchoices[answer=A]
{均匀变化的磁场能够在空间产生电场}
{电磁波在真空和介质中传播速度相同}
{只要有电场和磁场，就能产生电磁波}
{电磁波在同种介质中只能沿直线传播}
%% answer: A
%% memo:
\memoanswer{
变化的磁场可以产生电场，变化的电场可以产生磁场。A 正确。电磁波在真空和介质中传播速度不同，B 错误。电磁波是由变化的电场和变化的磁场由近及远的传播形成，C 错误。电磁波可以产生反射、折射、衍射和干涉，D 错误。
}

\item
%% number: 2
%% paperName: 2010年天津理综第 2 题 6 分
%% typeId: 1
%% type: 单选题
%% chapter: 原子和原子核
%% point: 玻尔模型
%% method: 无
%% score: 6
%% degree: 650
%% duplicateId: 0
%% body:
下列关于原子和原子核的说法正确的是\xzanswer{B}
\fourchoices[answer=B]
{$\beta$衰变现象说明电子是原子核的组成部分}
{玻尔理论的假设之一是原子能量的量子化}
{放射性元素的半衰期随温度的升高而变短}
{比结合能越小表示原子核中的核子结合得越牢固}
%% answer: B
%% memo:
\memoanswer{
$\beta$衰变是原子核中的质子转化为中子和电子，但电子不是原子核的组成部分，A 错误。玻尔理论的假设是把量子化代入到原子能量，B 正确。放射性元素的半衰期不随温度、状态及化学变化而变化，C 错误。比结合能越大表示原子核中的核子结合得越牢固，D 错误。
}

\item
%% number: 3
%% paperName: 2010年天津理综第 3 题 6 分
%% typeId: 1
%% type: 单选题
%% chapter: 直线运动
%% point: 运动图象
%% method: 无
%% score: 6
%% degree: 450
%% duplicateId: 0
%% body:
质点做直线运动的 $v-t$ 图像如图所示，规定向右为正方向，则该质点在前 $8\Us$ 内平均速度的大小和方向分别为\xzanswer{B}
\onepicture[align=right, w=5cm]{figs/03.png}
\fourchoices[answer=B]
{$0.25\Ums$，向右}
{$0.25\Ums$，向左}
{$1\Ums$，向右}
{$1\Ums$，向左}
%% answer: B
%% memo:
\memoanswer{
由题图得前 $8\Us$ 内的位移 $x=\frac{1}{2}\times3\times2\Um+\frac{1}{2}\times5\times(-2)\Um=-2\Um$，则平均速度 $\overline{v}=\frac{x}{t}=\frac{-2}{8}\Ums=-0.25\Ums$，负号表示方向向左。B 正确。
}

\item
%% number: 4
%% paperName: 2010年天津理综第 4 题 6 分
%% typeId: 1
%% type: 单选题
%% chapter: 机械波
%% point: 波的图象
%% method: 无
%% score: 6
%% degree: 650
%% duplicateId: 0
%% body:
一列简谐横波沿 $x$ 轴正向传播，传到 M 点时波形如图所示，再经 $0.6\Us$，N 点开始振动，则该波的振幅 $A$ 和频率 $f$ 为\xzanswer{D}
\onepicture[align=right, w=7cm]{figs/04.png}
\fourchoices[answer=D]
{$A=1\Um$，$f=5\UHz$}
{$A=0.5\Um$，$f=5\UHz$}
{$A=1\Um$，$f=2.5\UHz$}
{$A=0.5\Um$，$f=2.5\UHz$}
%% answer: D
%% memo:
\memoanswer{
由题图的振幅 $A=0.5\Um$，波长 $\lambda=4\Um$。经 $0.6\Us$，N 点开始振动得波速 $v=\frac{11-5}{0.6}\Ums=10\Ums$，所以频率 $f=\frac{v}{\lambda}=\frac{10}{4}\UHz=2.5\UHz$，D 正确。
}

\item
%% number: 5
%% paperName: 2010年天津理综第 5 题 6 分
%% typeId: 1
%% type: 单选题
%% chapter: 电场
%% point: 电势能电势差电场力做功
%% method: 无
%% score: 6
%% degree: 650
%% duplicateId: 0
%% body:
在静电场中，将一正电荷从 a 点移到 b 点，电场力做了负功，则\xzanswer{C}
\fourchoices[answer=C]
{b 点的电场强度一定比 a 点大}
{电场线方向一定从 b 指向 a}
{b 点的电势一定比 a 点高}
{该电荷的动能一定减小}
%% answer: C
%% memo:
\memoanswer{
正电荷从 a 移动到 b 点，电场力做了负功，电势能增加，说明 b 点的电势一定比 a 点高，选项 C 正确。电场强度的大小与电势的高低无关，无法确定，A 错误。b 点的电势比 a 点高，即 a、b 在不同的等势面上，但电场线方向不一定从 b 指向 a，B 错误。虽然电荷的电势能增加，如有重力做功，该电荷的动能不一定减小，D 错误。
}

\item
%% number: 6
%% paperName: 2010年天津理综第 6 题 6 分
%% typeId: 1
%% type: 单选题
%% chapter: 曲线运动
%% point: 人造地球卫星
%% method: 无
%% score: 6
%% degree: 650
%% duplicateId: 0
%% body:
探测器绕月球做匀速圆周运动，变轨后在周期较小的轨道上仍做匀速圆周运动，则变轨后与变轨前相比\xzanswer{A}
\fourchoices[answer=A]
{轨道半径变小}
{向心加速度变小}
{线速度变小}
{角速度变小}
%% answer: A
%% memo:
\memoanswer{
由于 $G\frac{Mm}{r^{2}}=m\left(\frac{2\pi}{T}\right)^{2}r$，所以 $r=\sqrt[3]{\frac{GMT^{2}}{4\pi^{2}}}$，$T$ 变小，$r$ 变小，A 正确。又 $G\frac{Mm}{r^{2}}=ma_{n}$，$a_{n}=\frac{GM}{r^{2}}$，$r$ 变小，$a_{n}$ 增大，B 错误。由 $G\frac{Mm}{r^{2}}=m\frac{v^{2}}{r}$，$v=\sqrt{\frac{GM}{r}}$，$r$ 变小，$v$ 增大，C 错误。由 $G\frac{Mm}{r^{2}}=m\omega^{2}r$，$\omega=\sqrt{\frac{GM}{r^{3}}}$，$r$ 变小，$\omega$ 增大，D 错误。
}

\item
%% number: 7
%% paperName: 2010年天津理综第 7 题 6 分
%% typeId: 2
%% type: 多选题
%% chapter: 交流电
%% point: 变压器
%% method: 无
%% score: 6
%% degree: 450
%% duplicateId: 0
%% body:
为探究理想变压器原、副线圈电压、电流的关系，将原线圈接到电压有效值不变的正弦交流电源上，副线圈连接相同的灯泡 L$_{1}$、L$_{2}$，电路中分别接了理想交流电压表 V$_{1}$、V$_{2}$ 和理想交流电流表 A$_{1}$、A$_{2}$，导线电阻不计，如图所示。当开关 S 闭合后\xzanswer{AD}
\onepicture[align=right, w=5.5cm]{figs/07.png}
\fourchoices[answer=AD]
{A$_{1}$ 示数变大，A$_{1}$ 与 A$_{2}$ 示数的比值不变}
{A$_{1}$ 示数变大，A$_{1}$ 与 A$_{2}$ 示数的比值变大}
{V$_{2}$ 示数变小，V$_{1}$ 与 V$_{2}$ 示数的比值变大}
{V$_{2}$ 示数不变，V$_{1}$ 与 V$_{2}$ 示数的比值不变}
%% answer: AD
%% memo:
\memoanswer{
由于理想变压器原线圈接到电压有效值不变，则副线圈电压不变，V$_{2}$ 示数不变，V$_{1}$ 与 V$_{2}$ 示数的比值不变，C 错误、D 正确。开关 S 闭合后，变压器副线圈的负载电阻减小，V$_{2}$ 不变，由欧姆定律可得 A$_{1}$ 示数变大，由于理想变压器 $P_{2}=P_{1}$，V$_{1}$ 与 V$_{2}$ 示数的比值不变，所以 A$_{1}$ 与 A$_{2}$ 示数的比值不变，A 正确、B 错误。正确选项 AD。
}

\item
%% number: 8
%% paperName: 2010年天津理综第 8 题 6 分
%% typeId: 2
%% type: 多选题
%% chapter: 光学
%% point: 光电效应
%% method: 无
%% score: 6
%% degree: 450
%% duplicateId: 0
%% body:
用同一光电管研究 a、b 两种单色光产生的光电效应，得到光电流 $I$ 与光电管两极间所加电压 $U$ 的关系如图。则这两种光\xzanswer{BC}
\onepicture[align=right, w=5cm]{figs/08.png}
\fourchoices[answer=BC]
{照射该光电管时 a 光使其逸出的光电子最大初动能大}
{从同种玻璃射入空气发生全反射时，a 光的临界角大}
{通过同一装置发生双缝干涉，a 光的相邻条纹间距大}
{通过同一玻璃三棱镜时，a 光的偏折程度大}
%% answer: BC
%% memo:
\memoanswer{
由光电效应方程 $\frac{1}{2}mv_{m}^{2}=h\nu-W_{0}$，由题图可得 b 光照射光电管时使其逸出的光电子最大初动能大，b 光的频率大，波长小。A 错误。b 光的频率大，在玻璃中的折射率 $n_{b}$ 大，由 $C=\arcsin\frac{1}{n}$ 从同种玻璃射入空气发生全反射时，b 光的临界角小，a 光大，B 正确。发生双缝干涉时，$\Delta x=\frac{L}{d}\lambda$，b 光波长小，相邻条纹间距 b 光小，a 光大，C 正确。在玻璃中的折射率 $n_{b}>n_{a}$，b 光的偏折程度大，D 错误。正确选项 BC。
}

\item
%% number: 9
%% paperName: 2010年天津理综第 9 题 4 分
%% typeId: 3
%% type: 填空题
%% chapter: 曲线运动
%% point: 平抛运动
%% method: 无
%% score: 4
%% degree: 450
%% duplicateId: 0
%% body:
如图所示，在高为 $h$ 的平台边缘水平抛出小球 A，同时在水平地面上距台面边缘水平距离为 $s$ 处竖直上抛小球 B，两球运动轨迹在同一竖直平面内，不计空气阻力，重力加速度为 $g$。若两球能在空中相遇，则小球 A 的初速度 $v_{A}$ 应大于\tkanswer[2.2cm]{$s\sqrt{\frac{g}{2h}}$}，A、B 两球初速度之比 $\frac{v_{A}}{v_{B}}$ 为\tkanswer[1.6cm]{$\frac{s}{h}$}。

\onepicture[align=right, w=5cm]{figs/09a.png}
%% answer: $s\sqrt{\frac{g}{2h}}$，$\frac{s}{h}$
%% memo:
\memoanswer{
由于 A 做平抛运动有 $s=v_{A}t$，$h=\frac{1}{2}gt^{2}$，要使两球能在空中相遇，所用时间小于 $t$，所以 $v_{A}>s\sqrt{\frac{g}{2h}}$。由 $s=v_{A}t'$，$h-h'=\frac{1}{2}gt'^{2}$，$h=v_{B}t'-\frac{1}{2}gt'^{2}$，得 $\frac{v_{A}}{v_{B}}=\frac{s}{h}$。
}

\item
%% number: 9
%% paperName: 2010年天津理综第 9 题 6 分
%% typeId: 4
%% type: 实验题
%% chapter: 力物体的平衡
%% point: 验证平行四边形实验
%% method: 无
%% score: 6
%% degree: 650
%% duplicateId: 0
%% body:
在探究合力的方法时，先将橡皮条的一端固定在水平木板上，另一端系上带有绳套的两根细绳。实验时，需要两次拉伸橡皮条，一次是通过两细绳用两个弹簧秤互成角度地拉橡皮条，另一次是用一个弹簧秤通过细绳拉橡皮条。

\begin{enumerate}
	\item
	实验对两次拉伸橡皮条的要求中，下列哪些说法是正确的\_\_\_\_\_\_\_\_（填字母代号）

	（A）将橡皮条拉伸相同长度即可

	（B）将橡皮条沿相同方向拉到相同长度

	（C）将弹簧秤都拉伸到相同刻度

	（D）将橡皮条和绳的结点拉到相同位置

	\item
	同学们在操作过程中有如下议论，其中对减小实验误差有益的说法是\_\_\_\_\_\_\_\_（填字母代号）

	（A）两细绳必须等长

	（B）弹簧秤、细绳、橡皮条都应与木板平行

	（C）用两弹簧秤同时拉细绳时两弹簧秤示数之差应尽可能大

	（D）拉橡皮条的细绳要长些，标记同一细绳方向的两点要远些
\end{enumerate}
%% answer: 
\jdanswer{
\begin{enumerate}
	\item
	BD

	\item
	BD

\end{enumerate}
}
%% memo:
\memoanswer{
（1）实验中两次拉伸橡皮条要将橡皮条拉伸相同长，即将橡皮条和绳的结点拉到相同位置。

（2）为了减小实验误差，尽可能使弹簧秤、细绳、橡皮条与木板平行；为了能更好地确定力的方向，细绳要长些，标记同一细绳方向的两点要远些。
}

\item
%% number: 9
%% paperName: 2010年天津理综第 9 题 8 分
%% typeId: 4
%% type: 实验题
%% chapter: 电路
%% point: 电路实验
%% method: 无
%% score: 8
%% degree: 450
%% duplicateId: 0
%% body:
要测量电压表 V$_{1}$ 的内阻 $R_{V}$，其量程为 $2\UV$，内阻约 $2\UkO$。实验室提供的器材有：

电流表 A，量程 $0.6\UA$，内阻约 $0.1\UO$；

电压表 V$_{2}$，量程 $5\UV$，内阻为 $5\UkO$；

定值电阻 $R_{1}$，阻值 $30\UO$；

定值电阻 $R_{2}$，阻值为 $3\UkO$；

滑动变阻器 $R_{3}$，最大阻值 $100\UO$，额定电流 $1.5\UA$；

电源 $E$，电动势 $6\UV$，内阻约 $0.5\UO$；

开关 S 一个，导线若干。

\begin{enumerate}
	\item
	有人拟将待测电压表 V$_{1}$ 和电流表 A 串联接入电压合适的测量电路中，测出 V$_{1}$ 的电压和电流，再计算出 $R_{V}$。该方案实际上不可行，其最主要的原因是\_\_\_\_\_\_\_\_\_

	\item
	请从上述器材中选择必要的器材，设计一个测量电压表 V$_{1}$ 内阻 $R_{V}$ 的实验电路。要求测量尽量准确，实验必须在同一电路中，且在不增减元件的条件下完成。试画出符合要求的实验电路图（图中电源与开关已连接好），并标出所选元件的相应字母代号：\_\_\_\_\_\_\_\_\_

	\item
	由上问写出 V$_{1}$ 内阻 $R_{V}$ 的表达方式，说明式中各测量量的物理意义。
\end{enumerate}

\onepicture[align=right, w=4cm]{figs/09b.png}
%% answer: 
\jdanswer{
\begin{enumerate}
	\item
	电流表 A 不能准确测量出流过电压表 V$_{1}$ 的电流

	\item
	实验电路如图所示

	\item
	$R_{V}=\frac{U_{1}R_{2}}{U_{2}-U_{1}}$，$U_{1}$ 表示 V$_{1}$ 的电压，$U_{2}$ 表示 V$_{1}$ 和 $R_{2}$ 串联的总电压

\end{enumerate}
}
%% memo:
\memoanswer{
（3）①由于电压表的内阻很大，流过的电流较小，待测电压表 V$_{1}$ 和电流表 A 串联接入电压合适的测量电路中，电流表 A 不能准确测量出流过电压表 V$_{1}$ 的电流。

②由于电压表 V$_{1}$ 接入电路后，由电流通过时其两端电压可以直接读出，因而利用串联分压的特点，选用标准电压表 V$_{2}$ 和定值电阻 $R_{2}$，反复测量通过滑动变阻器 $R_{3}$ 控制即可。测量电压表 V$_{1}$ 内阻 $R_{V}$ 的实验电路如图所示。

\onepicture[w=4cm]{figs/09_an.png}

③根据串联电路的特点：$\frac{U_{1}}{R_{V}}=\frac{U_{2}-U_{1}}{R_{2}}$（式中 $U_{1}$ 表示 V$_{1}$ 的电压，$U_{2}$ 表示 V$_{1}$ 和 $R_{2}$ 串联的总电压），所以 $R_{V}=\frac{U_{1}R_{2}}{U_{2}-U_{1}}$。
}

\item
%% number: 10
%% paperName: 2010年天津理综第 10 题 16 分
%% typeId: 6
%% type: 计算题
%% chapter: 运动和力
%% point: 动量和能量
%% method: 无
%% score: 16
%% degree: 450
%% duplicateId: 0
%% body:
如图所示，小球 A 系在细线的一端，线的另一端固定在 O 点，O 点到水平面的距离为 $h$。物块 B 质量是小球的 5 倍，置于粗糙的水平面上且位于 O 点的正下方，物块与水平面间的动摩擦因数为 $\mu$。现拉动小球使线水平伸直，小球由静止开始释放，运动到最低点时与物块发生正碰（碰撞时间极短），反弹后上升至最高点时到水平面的距离为 $\frac{h}{16}$。小球与物块均视为质点，不计空气阻力，重力加速度为 $g$，求物块在水平面上滑行的时间 $t$。

\onepicture[align=right, w=5.5cm]{figs/10.png}
%% answer: 
\jdanswer{
$\frac{\sqrt{2gh}}{4\mu g}$
}
%% memo:
\memoanswer{
设小球的质量为 $m$，运动到最低点与物块碰撞前的速度大小为 $v_{1}$，取小球运动到最低点重力势能为零，根据机械能守恒有：

$mgh=\frac{1}{2}mv_{1}^{2}$ ……①

得：$v_{1}=\sqrt{2gh}$。

设碰撞后小球反弹的速度大小为 $v_{1}'$，同理有：$mg\frac{h}{16}=\frac{1}{2}mv_{1}'^{2}$……②

得：$v_{1}'=\sqrt{\frac{gh}{8}}$。

设碰后物块的速度大小为 $v_{2}$，取水平向右为正方向，根据动量守恒定律有

$mv_{1}=-mv_{1}'+5mv_{2}$……③

得：$v_{2}=\sqrt{\frac{gh}{8}}$……④

物块在水平面上所受摩擦力的大小为：$F=5\mu mg$……⑤

设物块在水平面上滑行的时间为 $t$，根据动量定理有：$-Ft=0-5mv_{2}$……⑥

得：$t=\frac{\sqrt{2gh}}{4\mu g}$……⑦
}

\item
%% number: 11
%% paperName: 2010年天津理综第 11 题 18 分
%% typeId: 6
%% type: 计算题
%% chapter: 电磁感应
%% point: 电磁感应能量分析
%% method: 无
%% score: 18
%% degree: 450
%% duplicateId: 0
%% body:
如图所示，质量 $m_{1}=0.1\Ukg$，电阻 $R_{1}=0.3\UO$，长度 $l=0.4\Um$ 的导体棒 ab 横放在 U 型金属框架上。框架质量 $m_{2}=0.2\Ukg$，放在绝缘水平面上，与水平面间的动摩擦因数 $\mu=0.2$，相距 $0.4\Um$ 的 MM$'$、NN$'$ 相互平行，电阻不计且足够长。电阻 $R_{2}=0.1\UO$ 的 MN 垂直于 MM$'$。整个装置处于竖直向上的匀强磁场中，磁感应强度 $B=0.5\UT$。垂直于 ab 施加 $F=2\UN$ 的水平恒力，ab 从静止开始无摩擦地运动，始终与 MM$'$、NN$'$ 保持良好接触，当 ab 运动到某处时，框架开始运动。设框架与水平面间最大静摩擦力等于滑动摩擦力，$g$ 取 $10\Umsq$。

\begin{enumerate}
	\item
	求框架开始运动时 ab 速度 $v$ 的大小；
	\item
	从 ab 开始运动到框架开始运动的过程中，MN 上产生的热量 $Q=0.1\UJ$，求该过程 ab 位移 $x$ 的大小。
\end{enumerate}

\onepicture[align=right, w=6.5cm]{figs/11.png}
%% answer: 
\jdanswer{
\begin{enumerate}
	\item
	$v=6\Ums$

	\item
	$x=1.1\Um$

\end{enumerate}
}
%% memo:
\memoanswer{
（1）ab 对框架的压力 $F_{1}=m_{1}g$……①

框架受水平面的支持力 $F_{N}=m_{2}g+F_{1}$……②

依题意，最大静摩擦力等于滑动摩擦力，则框架受到最大静摩擦力 $F_{2}=\mu F_{N}$……③

ab 中的感应电动势 $E=Blv$……④

MN 中电流 $I=\frac{E}{R_{1}+R_{2}}$……⑤

MN 受到的安培力 $F_{\text{安}}=IlB$……⑥

框架开始运动时 $F_{\text{安}}=F_{2}$……⑦

由上述各式代入数据解得 $v=6\Ums$……⑧

（2）闭合回路中产生的总热量：$Q_{\text{总}}=\frac{R_{1}+R_{2}}{R_{2}}Q$……⑨

由能量守恒定律，得：$Fx=\frac{1}{2}m_{1}v^{2}+Q_{\text{总}}$……⑩

代入数据解得 $x=1.1\Um$……(11)
}

\item
%% number: 12
%% paperName: 2010年天津理综第 12 题 20 分
%% typeId: 6
%% type: 计算题
%% chapter: 磁场
%% point: 带电粒子在复合场中的运动
%% method: 无
%% score: 20
%% degree: 250
%% duplicateId: 0
%% body:
质谱分析技术已广泛应用于各前沿科学领域。汤姆孙发现电子的质谱装置示意如图，M、N 为两块水平放置的平行金属极板，板长为 $L$，板右端到屏的距离为 $D$，且 $D$ 远大于 $L$，O$'$O 为垂直于屏的中心轴线，不计离子重力和离子在板间偏离 O$'$O 的距离。以屏中心 O 为原点建立 $x$O$y$ 直角坐标系，其中 $x$ 轴沿水平方向，$y$ 轴沿竖直方向。

上述装置中，保留原电场，再在板间加沿 $-y$ 方向的匀强磁场。现有电荷量相同的两种正离子组成的离子流，仍从 O$'$ 点沿 O$'$O 方向射入，屏上出现两条亮线。在两线上取 $y$ 坐标相同的两个光点，对应的 $x$ 坐标分别为 $3.24\Umm$ 和 $3.00\Umm$，其中 $x$ 坐标大的光点是碳 12 离子击中屏产生的，另一光点是未知离子产生的。尽管入射离子速度不完全相同，但入射速度都很大，且在板间运动时 O$'$O 方向的分速度总是远大于 $x$ 方向和 $y$ 方向的分速度。

\begin{enumerate}
	\item
	设一个质量为 $m_{0}$、电荷量为 $q_{0}$ 的正离子以速度 $v_{0}$ 沿 O$'$O 的方向从 O$'$ 点射入，板间不加电场和磁场时，离子打在屏上 O 点。若在两极板间加一沿 $+y$ 方向场强为 $E$ 的匀强电场，求离子射到屏上时偏离 O 点的距离 $y_{0}$；
	\item
	假设你利用该装置探究未知离子，试依照以下实验结果计算未知离子的质量数。
\end{enumerate}

\onepicture[align=right, w=7.5cm]{figs/12.png}
%% answer: 
\jdanswer{
\begin{enumerate}
	\item
	$\frac{q_{0}ELD}{m_{0}v_{0}^{2}}$

	\item
	14

\end{enumerate}
}
%% memo:
\memoanswer{
（1）离子在电场中受到的电场力 $F_{y}=q_{0}E$……①

离子获得的加速度 $a_{y}=\frac{F_{y}}{m_{0}}$……②

离子在板间运动的时间 $t_{0}=\frac{L}{v_{0}}$……③

到达极板右边缘时，离子在 $+y$ 方向的分速度 $v_{y}=a_{y}t_{0}$……④

离子从板右端到达屏上所需时间 $t_{0}'=\frac{D}{v_{0}}$……⑤

离子射到屏上时偏离 O 点的距离 $y_{0}=v_{y}t_{0}'$

由上述各式，得 $y_{0}=\frac{q_{0}ELD}{m_{0}v_{0}^{2}}$……⑥

（2）设离子电荷量为 $q$，质量为 $m$，入射时速度为 $v$，磁场的磁感应强度为 $B$，磁场对离子的洛伦兹力 $F_{x}=qvB$……⑦

已知离子的入射速度都很大，因而离子在磁场中运动时间甚短。所经过的圆弧与圆周相比甚小，且在板间运动时 O$'$O 方向的分速度总是远大于在 $x$ 方向和 $y$ 方向的分速度，洛伦兹力变化甚微，故可作恒力处理，洛伦兹力产生的加速度 $a_{x}=\frac{qvB}{m}$……⑧

$a_{x}$ 是离子在 $x$ 方向的加速度，离子在 $x$ 方向的运动可视为初速度为零的匀加速直线运动，到达极板右端时，离子在 $x$ 方向的分速度 $v_{x}=a_{x}t=\frac{qvB}{m}\left(\frac{L}{v}\right)=\frac{qBL}{m}$……⑨

离子飞出极板到达屏时，在 $x$ 方向上偏离 O 点的距离 $x=v_{x}t'=\frac{qBL}{m}\left(\frac{D}{v}\right)=\frac{qBLD}{mv}$……⑩

当离子的初速度为任意值时，离子到达屏上时的位置在 $y$ 方向上偏离 O 点的距离为 $y$，考虑到⑥式，得 $y=\frac{qELD}{mv^{2}}$……(11)

由⑩、(11)两式得 $x^{2}=\frac{k}{m}y$……(12)

其中 $k=\frac{qB^{2}LD}{E}$

上式表明，$k$ 是与离子进入板间初速度无关的定值，对两种离子均相同，由题设条件知，$x$ 坐标 $3.24\Umm$ 的光点对应的是碳 12 离子，其质量为 $m_{1}=12\Uu$，$x$ 坐标 $3.00\Umm$ 的光点对应的是未知离子，设其质量为 $m_{2}$，由(12)式代入数据可得 $m_{2}\approx14\Uu$……(13)

故该未知离子的质量数为 14。
}

\end{enumerate}

\end{document}
